\section{Metodologia}\label{sec:km_metodo}

\textbf{Pré-processamento.} Usamos o mesmo pré-processamento da parte anterior do trabalho, em que treinamos máquinas de vetores de suporte (SVM)~\cite{cortes:95} para prever a renda: os 5 atributos numéricos padronizados e os 5 de texto em one-hot, num total de \clColunas{} colunas. A SVM com kernel RBF também mede distâncias, e o mesmo cuidado com a escala vale aqui. A única diferença é um logaritmo, $\log(1 + x)$, no ganho e na perda de capital antes da padronização; o ``1 +'' existe porque a maioria tem valor zero, e $\log 0$ não existe. A Seção~\ref{sec:cl_log} compara os grupos com e sem esse logaritmo.

\textbf{Protocolo.} Não há divisão em treino e teste, porque não há resposta a acertar: os grupos são formados na base inteira, e não há vazamento, porque a renda nunca entra no ajuste. Para cada $k$ de 2 a 10, o K-means roda com 10 seeds (0 a 9), e cada seed faz 10 inicializações com k-means++ e fica com a de menor inércia. A seed fixa o gerador de números aleatórios, e portanto quais centros iniciais são sorteados. Em cada ajuste, medimos a inércia, a silhueta e o índice de Davies-Bouldin. A silhueta compara todos os pares de pessoas e é calculada numa amostra fixa de 10 mil pessoas, a mesma em todos os ajustes. A estabilidade de cada $k$ é o ARI médio entre as partições das 10 seeds, comparadas duas a duas (45 pares), e repetimos o estudo com uma só inicialização por seed, para medir o que as 10 inicializações acrescentam.

\textbf{Partição descrita.} Escolhido o $k$, descrevemos a partição da seed 42, fora das seeds 0 a 9, como nas análises de uma seed das partes anteriores do trabalho. Não faz sentido tirar a média das 10 partições: o número que o K-means dá a cada grupo é arbitrário e muda de uma seed para outra. Para facilitar a leitura, os grupos foram renumerados de 1 a $k$ pelo percentual de renda alta; a renda só ordena os nomes, sem influenciar a formação dos grupos.

\textbf{Implementação.} Tudo foi feito com o scikit-learn~\cite{pedregosa:11}, e os tempos foram medidos num único núcleo de processador. O código, os resultados e o notebook com a análise completa estão no repositório do trabalho, na pasta \texttt{renda/}.
